\documentclass[11pt]{article}
\usepackage[margin=1in]{geometry}
\usepackage{amsmath,amssymb,amsthm}
\usepackage{xurl}
\usepackage{hyperref}
\sloppy
\let\oldus\_ \renewcommand{\_}{\oldus\allowbreak}
\newtheorem{theorem}{Theorem}
\newtheorem{lemma}[theorem]{Lemma}
\newtheorem{corollary}[theorem]{Corollary}
\theoremstyle{definition}
\newtheorem{definition}[theorem]{Definition}
\newcommand{\vol}{\operatorname{vol}}
\newcommand{\ip}[2]{\langle #1,#2\rangle}

\title{Disproof of conjecture 00000007792:\\ the M-ellipsoid constant of the Euclidean ball is $4$, not $\le (\pi/4)e$}
\author{Jackmeson1}
\date{2026-10-05}

\begin{document}
\maketitle

\section{The conjecture and the clause refuted}

\textbf{Statement (English, verbatim).} Definition: The M-ellipsoid condition: an M-position of a
convex body $K$ is an ellipsoid $E$ with
$\vol(K+E)\,\vol(K^\circ+E^\circ)\le C^n\,\vol(K)\,\vol(K^\circ)$. Conjecture: The optimal constant
$C$ is uniformly attainable over all bodies with $C\le(\pi/4)e$, given asymptotically by the John
ellipsoid position; and M-positions are stable: bodies deviating from the optimal position lose
exponentially and strictly monotonically in volume product, with characterization function
$e^{c d^2}$ in the position deviation $d$. (The Chinese text states the same: the optimal constant
$C$ is uniformly attainable over all bodies and $C\le \pi/4\cdot e$.)

\textbf{Reading.} The conjecture is a conjunction. Its first conjunct says that a single constant
$C\le(\pi/4)e\approx 2.135$ works for all convex bodies $K$ (in every dimension $n$), i.e.\ every
$K$ has an ellipsoid $E$ with
\[
  \vol(K+E)\,\vol(K^\circ+E^\circ)\le C^n\,\vol(K)\,\vol(K^\circ),
\]
the inequality exactly as written in the Definition line (with $C^n$, not $C^{2n}$). We refute this
conjunct, which refutes the conjunction. The stability conjunct (the $e^{cd^2}$ law) is not
addressed. We refute three readings of the first conjunct:
\begin{itemize}
\item[(R1)] one $C\le(\pi/4)e$ is admissible in every dimension $n\ge 1$;
\item[(R2)] (``asymptotically'') one $C\le(\pi/4)e$ is admissible in all sufficiently large
  dimensions;
\item[(R3)] admissible constants $C_n\ge 0$ in dimension $n$ (for all large $n$) converge to a
  limit $L\le(\pi/4)e$.
\end{itemize}
In fact, in every dimension $n\ge1$ every admissible constant satisfies $|C|\ge 4$. A reading with
$C\le \pi/(4e)$ is stronger than (R1) and is refuted too. The witness is the Euclidean unit ball,
and the bound $4$ holds against \emph{every} ellipsoid,
so no choice of position (John or otherwise) helps.

\textbf{Conventions.} Space $\mathbb{R}^n$ with the Euclidean inner product (Mathlib
\texttt{EuclideanSpace R (Fin n)}); $\vol$ is Lebesgue measure (Mathlib \texttt{volume}, valued in
$[0,\infty]$); $K+E$ is the Minkowski sum. Polarity is with respect to the origin. Ellipsoids may
have any centre. The constant $C$ is real; we prove $|C|\ge4$ in each dimension, and $C\ge 4$ in
odd dimensions (where $C^n<0$ is excluded automatically).

\section{Definitions (as formalized)}

\begin{definition}
$B=\{x:\|x\|\le1\}$. The \emph{polar} of $K\subseteq\mathbb{R}^n$ is
$K^\circ=\{y:\ \ip{x}{y}\le1\ \text{for all } x\in K\}$ (\texttt{polar}). An \emph{ellipsoid} is a
set $E=\{c+Ax:\ x\in B\}$ with $c\in\mathbb{R}^n$ and $A$ an invertible linear map
(\texttt{IsEllipsoid}). A \emph{convex body} is a compact convex set with nonempty interior
(\texttt{IsConvexBody}).
\end{definition}

\begin{definition}
$E$ satisfies the \emph{M-ellipsoid condition} for $K$ with constant $C$
(\texttt{MEllipsoidCondition}) if $E$ is an ellipsoid and
$\vol(K+E)\vol(K^\circ+E^\circ)\le C^n\vol(K)\vol(K^\circ)$, where $C^n$ enters through
$\max(C^n,0)$ (\texttt{ENNReal.ofReal}). $C$ is \emph{admissible in dimension $n$}
(\texttt{AdmissibleConstant n C}) if every convex body $K\subseteq\mathbb{R}^n$ with
$0\in\operatorname{int}K$ has an ellipsoid $E$ satisfying the condition.
\end{definition}

\section{Proof}

\begin{lemma}[\texttt{polar\_closedBall}, \texttt{closedBall\_isConvexBody}]
$B^\circ=B$, and $B$ is a convex body with $0$ in its interior.
\end{lemma}
\begin{proof}
If $\|y\|\le 1$ then $\ip{x}{y}\le\|x\|\|y\|\le1$ on $B$. If $y\in B^\circ$, $y\ne0$, test
$x=y/\|y\|$ to get $\|y\|\le 1$.
\end{proof}

\begin{lemma}[\texttt{key\_ineq}]
For $d,x\in\mathbb{R}^n$ with $\|x\|\le1$ and $k=(1+\|d\|)^{-1}$:
$\ip{d}{x-kd}+\|x-kd\|\le1$.
\end{lemma}
\begin{proof}
Put $r=\|d\|$, $p=\ip{d}{x}\le r$. Then $\ip{d}{x-kd}=p-kr^2$, and $R:=1-p+kr^2\ge 1-r+kr^2=k>0$.
Using $k(1+r)=1$ one checks the identity
$R^2-\|x-kd\|^2=(r-p)^2+1-\|x\|^2\ge0$, so $\|x-kd\|\le R$.
\end{proof}

\begin{theorem}[\texttt{ball\_volume\_product\_ge}]
For every $n$ and every ellipsoid $E\subseteq\mathbb{R}^n$:
$4^n\vol(B)\vol(B^\circ)\le\vol(B+E)\,\vol(B^\circ+E^\circ)$.
\end{theorem}
\begin{proof}
Let $E=c+AB$. The map $A^*A$ is symmetric; let $v_1,\dots,v_n$ be an orthonormal eigenbasis,
$A^*Av_j=\lambda_jv_j$ (Mathlib \texttt{LinearMap.IsSymmetric.eigenvectorBasis}). Then
$\ip{Av_i}{Av_j}=\lambda_j\delta_{ij}$, so $\lambda_j=\|Av_j\|^2>0$. Put $s_j=\sqrt{\lambda_j}$ and
$u_j=s_j^{-1}Av_j$; the $u_j$ form an orthonormal basis. Let $R$ be the isometry with $Ru_j=v_j$,
and for $a\in\mathbb{R}^n$ let $D_a$ be the linear map with $D_au_j=a_ju_j$, so
$\det D_a=\prod_j a_j$ and $D_{1+a}=I+D_a$. Checking on the basis $u_j$:
$A R=D_s$ and $A^*D_{s^{-1}}=R$.

\emph{First factor.} For $x\in B$, $c+D_{1+s}x=x+(c+A(Rx))\in B+E$ because $\|Rx\|=\|x\|$. Hence
$B+E\supseteq c+D_{1+s}B$ and, by translation invariance and the linear change of variables
(\texttt{addHaar\_image\_linearMap}), $\vol(B+E)\ge\prod_j(1+s_j)\,\vol(B)$.

\emph{Second factor.} Let $d=D_{s^{-1}}^*c$ and $k=(1+\|d\|)^{-1}$. For $x,x'\in B$ and
$z=x-kd$:
\[
 \ip{c+Ax'}{D_{s^{-1}}z}=\ip{d}{z}+\ip{x'}{A^*D_{s^{-1}}z}=\ip{d}{z}+\ip{x'}{Rz}
 \le\ip{d}{z}+\|z\|\le1
\]
by the key inequality. So $D_{s^{-1}}(x-kd)\in E^\circ$, and
$x+D_{s^{-1}}(x-kd)=-kD_{s^{-1}}d+D_{1+s^{-1}}x\in B+E^\circ$. Hence
$\vol(B^\circ+E^\circ)=\vol(B+E^\circ)\ge\prod_j(1+s_j^{-1})\,\vol(B)$.

\emph{Product.} $(1+s)(1+s^{-1})-4=(s-1)^2/s\ge0$ for $s>0$
(\texttt{four\_le\_mul\_inv}), so the product of the two lower bounds is at least
$4^n\vol(B)^2=4^n\vol(B)\vol(B^\circ)$.
\end{proof}

\begin{corollary}[\texttt{ball\_condition\_pow\_ge}, \texttt{admissible\_abs\_ge\_four},
\texttt{admissible\_ge\_four\_of\_odd}]
If some ellipsoid satisfies the M-ellipsoid condition for $B$ with constant $C$, then
$4^n\le C^n$. Hence in every dimension $n\ge1$ every admissible $C$ has $|C|\ge4$, and $C\ge4$ when
$n$ is odd.
\end{corollary}
\begin{proof}
Combine the theorem with the hypothesis and cancel $\vol(B)^2\in(0,\infty)$; if $C^n\le0$ the
right side is $0$, contradicting $4^n>0$.
\end{proof}

\begin{corollary}[\texttt{uniform\_admissible\_ge\_four}, \texttt{conjecture\_7792\_false},
\texttt{conjecture\_7792\_false\_eventually}, \texttt{conjecture\_7792\_false\_limit}]
Since $(\pi/4)e<e<3<4$: (R1) fails (a uniformly admissible $C$ is $\ge4$, using $n=1$); (R2) fails
(use the odd dimension $2N+1$); (R3) fails (eventually $C_n\ge4$, so $L\ge4$).
\end{corollary}

The bound is sharp for the ball: $E=B$ gives $\vol(2B)^2=4^n\vol(B)^2$ (not formalized; not needed).

\section{Lean formalization and axiom audit}

File \texttt{lean/Conjecture7792/Basic.lean} (332 lines, \texttt{import Mathlib} only, Lean
\texttt{v4.33.1}, Mathlib \texttt{v4.33.1}). Main theorem:
\begin{verbatim}
theorem conjecture_7792_false :
    Not (exists C : Real, C <= Real.pi / 4 * Real.exp 1 /\
      forall n : Nat, 1 <= n -> AdmissibleConstant n C)
\end{verbatim}
(ASCII rendering of the Lean statement)
together with \texttt{conjecture\_7792\_false\_eventually},
\texttt{conjecture\_7792\_false\_limit}, \texttt{admissible\_abs\_ge\_four} and the core estimate
\texttt{ball\_volume\_product\_ge}. \texttt{lake build} succeeds. \texttt{\#print axioms} for each
of these reports only \texttt{propext}, \texttt{Classical.choice}, \texttt{Quot.sound}. There is
no \texttt{sorry} and no \texttt{native\_decide}.

\end{document}
